\documentclass[article]{jss}\usepackage[]{graphicx}\usepackage[]{color}
% maxwidth is the original width if it is less than linewidth
% otherwise use linewidth (to make sure the graphics do not exceed the margin)
\makeatletter
\def\maxwidth{ %
  \ifdim\Gin@nat@width>\linewidth
    \linewidth
  \else
    \Gin@nat@width
  \fi
}
\makeatother

\definecolor{fgcolor}{rgb}{0.345, 0.345, 0.345}
\newcommand{\hlnum}[1]{\textcolor[rgb]{0.686,0.059,0.569}{#1}}%
\newcommand{\hlstr}[1]{\textcolor[rgb]{0.192,0.494,0.8}{#1}}%
\newcommand{\hlcom}[1]{\textcolor[rgb]{0.678,0.584,0.686}{\textit{#1}}}%
\newcommand{\hlopt}[1]{\textcolor[rgb]{0,0,0}{#1}}%
\newcommand{\hlstd}[1]{\textcolor[rgb]{0.345,0.345,0.345}{#1}}%
\newcommand{\hlkwa}[1]{\textcolor[rgb]{0.161,0.373,0.58}{\textbf{#1}}}%
\newcommand{\hlkwb}[1]{\textcolor[rgb]{0.69,0.353,0.396}{#1}}%
\newcommand{\hlkwc}[1]{\textcolor[rgb]{0.333,0.667,0.333}{#1}}%
\newcommand{\hlkwd}[1]{\textcolor[rgb]{0.737,0.353,0.396}{\textbf{#1}}}%
\let\hlipl\hlkwb

\usepackage{framed}
\makeatletter
\newenvironment{kframe}{%
 \def\at@end@of@kframe{}%
 \ifinner\ifhmode%
  \def\at@end@of@kframe{\end{minipage}}%
  \begin{minipage}{\columnwidth}%
 \fi\fi%
 \def\FrameCommand##1{\hskip\@totalleftmargin \hskip-\fboxsep
 \colorbox{shadecolor}{##1}\hskip-\fboxsep
     % There is no \\@totalrightmargin, so:
     \hskip-\linewidth \hskip-\@totalleftmargin \hskip\columnwidth}%
 \MakeFramed {\advance\hsize-\width
   \@totalleftmargin\z@ \linewidth\hsize
   \@setminipage}}%
 {\par\unskip\endMakeFramed%
 \at@end@of@kframe}
\makeatother

\definecolor{shadecolor}{rgb}{.97, .97, .97}
\definecolor{messagecolor}{rgb}{0, 0, 0}
\definecolor{warningcolor}{rgb}{1, 0, 1}
\definecolor{errorcolor}{rgb}{1, 0, 0}
\newenvironment{knitrout}{}{} % an empty environment to be redefined in TeX

\usepackage{alltt}

%% -- LaTeX packages and custom commands ---------------------------------------

%% recommended packages
\usepackage{thumbpdf,lmodern}
\usepackage{float,tikz}
%% another package (only for this demo article)
\usepackage{framed}
%% author pakckage
\usepackage{xcolor,caption, float}

\graphicspath{ {./imgs/} }

%% new custom commands
\newcommand{\class}[1]{`\code{#1}'}
\newcommand{\fct}[1]{\code{#1()}}

% <<preliminaries, echo=FALSE, results=hide>>=


%% -- Article metainformation (author, title, ...) -----------------------------

%% - \author{} with primary affiliation
%% - \Plainauthor{} without affiliations
%% - Separate authors by \And or \AND (in \author) or by comma (in \Plainauthor).
%% - \AND starts a new line, \And does not.
\author{Thomas Huet\\UMR 5140}
%   \And Second Author\\Plus Affiliation}
\Plainauthor{Thomas Huet}

%% - \title{} in title case
%% - \Plaintitle{} without LaTeX markup (if any)
%% - \Shorttitle{} with LaTeX markup (if any), used as running title
\title{Modelling Prehistorical Iconographic Compositions. The \proglang{R} package \pkg{decorr}}
\Plaintitle{Modelling Prehistorical Iconographic Compositions. The R package 'decorr'}
\Shorttitle{Modelling Prehistorical Iconographic Compositions}

%% - \Abstract{} almost as usual
\Abstract{
  By definition, Prehistorical societies are characterised by the absence of a writing system. Prehistorical times cover more than 99\% of the human living. Even if it is being discussed, first symbolic manifestations start around 200,000 BC \citep{dErrico00}. The duration from first symbolic expressions to start of writing represents 97\% of the human living. In illiterate societies, testimonies of symbolic systems mostly come from iconography (ceramic decorations, rock-art, statuary, etc.) and signs are displayed mostlty a discontinuous figures which can have different relationships one with another. An graphical composition can be "read" as a spatial distribution of features having intrinsic values possibily having meaningful relationships one with another depending on their pairwise spatial proximities. 
}

\Keywords{Iconography, Prehistory, Graph Theory, Graph Drawing, Spatial Analysis, \proglang{R}}
\Plainkeywords{Iconography, Prehistory, Graph Theory, Graph Drawing, Spatial Analysis, R}

\Address{
  Thomas Huet\\
  CNRS-UMR 5140 \\
  Archeologie des Societes Mediterraneennes\\
  Universite Paul Valery\\
  route de Mende\\
  Montpellier 34199, France\\
  E-mail: \email{thomashuet7@gmail.com}\\
%  URL: \url{https://eeecon.uibk.ac.at/~zeileis/}
}
\IfFileExists{upquote.sty}{\usepackage{upquote}}{}
\begin{document}
% \Sweave2knitr("article_rvTH14_1.Rnw")
\SweaveOpts{concordance=TRUE}

\section[Introduction]{Introduction} \label{sec:intro}

Until our days, formal methods to study ancient iconography Semantics, has been mostlty been grounded (explicitly or not) on the prime principle of Saussurian linguistic: the 'linearity of the signifier' \citep{Saussure89}.  
Writing is one of the most rational semiographical system. With a clear distinction between signified and signifier -- specially in alphabetic and binary writings -- and the development of the signified on a horizontal, vertical or boustrophedon axis. 
Let us take the example of the word "\code{art}" which contains three vertices (\code{a}, \code{r}, \code{t}) and two edges (one between \code{a} and \code{r}, the other between \code{r} and \code{t}). In \proglang{R}, these features, concatenated in this order with a \code{paste0()}, is \code{art}, and not \code{rat}

% \SweaveOpts{width=3,height=1}
\begin{figure}[H]
\begin{knitrout}
\definecolor{shadecolor}{rgb}{0.969, 0.969, 0.969}\color{fgcolor}\begin{kframe}
\begin{alltt}
\hlstd{R> }\hlkwd{library}\hlstd{(igraph);}\hlkwd{library}\hlstd{(extrafont)}
\hlstd{R> }\hlkwd{par}\hlstd{(}\hlkwc{mar}\hlstd{=}\hlkwd{c}\hlstd{(}\hlnum{0.1}\hlstd{,}\hlnum{0.1}\hlstd{,}\hlnum{0.1}\hlstd{,}\hlnum{0.1}\hlstd{) )}
\hlstd{R> }\hlcom{# loadfonts(device = "win")}
\hlstd{R> }\hlcom{# ugs <- data.frame(name=c("a","r","t"))}
\hlstd{R> }\hlcom{# lks <- data.frame(from=c("a", "r"),to=c("r", "t"),type=c("=","="))}
\hlstd{R> }\hlcom{# g <- graph_from_data_frame(lks, directed=FALSE, vertices= ugs)}
\hlstd{R> }\hlstd{g} \hlkwb{<-} \hlkwd{graph_from_literal}\hlstd{(a}\hlopt{-+}\hlstd{r, r}\hlopt{-+}\hlstd{t)}
\hlstd{R> }\hlcom{# lo <- cbind(seq(-1,1,length.out = gorder(g)), 0)}
\hlstd{R> }\hlkwd{plot}\hlstd{(g,}
\hlstd{+  }     \hlkwc{layout} \hlstd{=} \hlkwd{cbind}\hlstd{(}\hlkwd{seq}\hlstd{(}\hlopt{-}\hlnum{1}\hlstd{,}\hlnum{1}\hlstd{,}\hlkwc{length.out} \hlstd{=} \hlkwd{gorder}\hlstd{(g)),} \hlnum{0}\hlstd{),}
\hlstd{+  }     \hlkwc{vertex.color}\hlstd{=}\hlstr{"white"}\hlstd{,}
\hlstd{+  }     \hlkwc{vertex.frame.color}\hlstd{=}\hlstr{"white"}\hlstd{,}
\hlstd{+  }     \hlkwc{vertex.size}\hlstd{=}\hlnum{60}\hlstd{,}
\hlstd{+  }     \hlkwc{vertex.label.cex}\hlstd{=}\hlnum{1.5}\hlstd{,}
\hlstd{+  }     \hlkwc{vertex.label.color}\hlstd{=}\hlstr{"black"}\hlstd{,}
\hlstd{+  }     \hlkwc{vertex.label.family}\hlstd{=}\hlstr{"Courier New"}\hlstd{,}
\hlstd{+  }     \hlkwc{edge.arrow.size}\hlstd{=}\hlnum{.5}\hlstd{)}
\end{alltt}
\end{kframe}
\includegraphics[width=\maxwidth]{figure/unnamed-chunk-1-1} 

\end{knitrout}
\captionof{figure}[directional links]{concatenate of \code{a}, \code{r} and \code{t} graphical units (GUs) is \code{art}.}
\end{figure}

But, as stated, in Prehistorical the writing system does not exists. Spatial relationships between graphical features, or graphical units (GUs) are not necessarly linear and directed but could most probably be more multi-directional and undirected: the direction of the interactions of pairwise GUs can be in any order.  

\end{document}
